\documentclass[../../main.tex]{subfiles}
\graphicspath{{../../image/}{../../figure/}}
\begin{document}

\section{特征线的约束与锯齿剖面的形成}\label{section:ch11-特征线约束与锯齿剖面}

极值后向特征线在 $x$ 轴有界区间上的截距受到约束, 这诱导出特征线速度递减变差的界, 从而由真正非线性, 诱导出解本身递减变差的界。

\begin{theorem}\label{theorem:ch11-11-7-1}
设 $\chi_-(\cdot)$ 与 $\chi_+(\cdot)$ 是 $[0,\infty)$ 上与可容许解 $u$ 关联的广义特征线, 且对 $t\in[0,\infty)$, $\chi_-(t)<\chi_+(t)$。则对任意 $t>0$, 函数 $f'(u(\cdot,t))$ 在区间 $(\chi_-(t),\chi_+(t))$ 上的递减变差不超过 $[\chi_+(0)-\chi_-(0)]t^{-1}$。因而 $u(\cdot,t)$ 在区间 $(\chi_-(t),\chi_+(t))$ 上的递减变差当 $t\to\infty$ 时为 $O(t^{-1})$。
\end{theorem}

\begin{proof}
固定 $t>0$, 考虑任意分划 $\chi_-(t)<x_1<x_2<\cdots<x_{2m}<\chi_+(t)$, 使得 $(x_i,t)$ 是 $u$ 的连续点, 且 $u(x_{2k-1},t)>u(x_{2k},t)$, $k=1,\cdots,m$。设 $\xi_i(\cdot)$ 表示从 $(x_i,t)$ 出发的(唯一)后向特征线。则 $\chi_-(0)\leqslant\xi_1(0)\leqslant\cdots\leqslant\xi_{2m}(0)\leqslant\chi_+(0)$。此外, $\xi_i(0)=x_i-tf'(u(x_i,t))$, 于是
\begin{align}
\sum_{k=1}^m t[f'(u(x_{2k-1},t))-f'(u(x_{2k},t))]\leqslant\chi_+(0)-\chi_-(0),\label{align:ch11-11-7-1}
\end{align}
由此得定理断言。证毕。
\end{proof}

特别地, 参照\cref{theorem:ch11-11-6-1}的设定, 我们推出 $N$ 波剖面 $u(\cdot,t)$ 在区间 $(\chi_-(t),\chi_+(t))$ 上的递减变差当 $t\to\infty$ 时为 $O(t^{-1})$。

\cref{theorem:ch11-11-7-1}的另一推论是: 当初值 $u_0$, 从而解 $u$, 关于 $x$ 周期时, 任意周期区间上 $u(\cdot,t)$ 的递减变差, 从而也全变差, 当 $t\to\infty$ 时为 $O(t^{-1})$。通过更细致地考察初值, 我们可以获得比 $O(t^{-1})$ 更精细的分辨率。

\begin{theorem}\label{theorem:ch11-11-7-2}
设 $u$ 是初值为 $u_0$ 的可容许解。假设 $\chi_-(t)=x_-+tf'(\bar u)$ 与 $\chi_+(t)=x_++tf'(\bar u)$, $x_-<x_+$, 是与 $u$ 关联的相邻分界线, 即\eqref{align:ch11-11-4-2}对 $\bar x=x_-$ 和 $\bar x=x_+$ 成立, 但对区间 $(x_-,x_+)$ 内任何 $\bar x$ 不成立。则
\begin{align}
\int_{\chi_-(t)}^{\chi_+(t)}u(x,t)\,\mathrm{d}x=\int_{x_-}^{x_+}u_0(y)\,\mathrm{d}y=(x_+-x_-)\bar u,\quad t\in[0,\infty).\label{align:ch11-11-7-2}
\end{align}
考虑从点 $\left(\frac{x_-+x_+}{2},0\right)$ 出发的任意前向特征线 $\psi(\cdot)$。则当 $t\to\infty$ 时,
\begin{align}
\psi(t)=\tfrac12[\chi_-(t)+\chi_+(t)]+o(1),\label{align:ch11-11-7-3}
\end{align}
\begin{align}
u(x\pm,t)=
\begin{cases}
\bar u+\dfrac{1}{f''(\bar u)}\dfrac{x-\chi_-(t)}{t}+o\left(\dfrac1t\right), & \chi_-(t)<x<\psi(t),\\[3mm]
\bar u+\dfrac{1}{f''(\bar u)}\dfrac{x-\chi_+(t)}{t}+o\left(\dfrac1t\right), & \psi(t)<x<\chi_+(t).
\end{cases}\label{align:ch11-11-7-4}
\end{align}
此外, $u(\cdot,t)$ 在区间 $(\chi_-(t),\psi(t))$ 与 $(\psi(t),\chi_+(t))$ 上的递减变差当 $t\to\infty$ 时为 $o(t^{-1})$。
\end{theorem}

\begin{proof}
为验证\eqref{align:ch11-11-7-2}中第一等式, 只需在平行四边形 $\{(x,\tau):0<\tau<t,\ \chi_-(\tau)<x<\chi_+(\tau)\}$ 上积分\eqref{align:ch11-11-1-1}, 再应用 Green 公式。\eqref{align:ch11-11-7-2}中第二等式成立, 因为\eqref{align:ch11-11-4-2}对 $\bar x=x_-$ 与 $\bar x=x_+$ 均成立。

对 $t>0$, 设 $\xi_-^t(\cdot)$ 与 $\xi_+^t(\cdot)$ 表示从点 $(\psi(t),t)$ 出发的极值后向特征线。当 $t\uparrow\infty$ 时, $\xi_-^t(0)\downarrow x_-$ 且 $\xi_+^t(0)\uparrow x_+$, 否则将存在分界线发自点 $(\bar x,0)$, $\bar x\in(x_-,x_+)$, 与假设矛盾。于是由\cref{theorem:ch11-11-7-1}, $f'(u(\cdot,t))$ 的递减变差, 从而 $u(\cdot,t)$ 本身在区间 $(\chi_-(t),\psi(t))$ 与 $(\psi(t),\chi_+(t))$ 上的递减变差, 当 $t\to\infty$ 时为 $o(t^{-1})$。

从任意点 $(x,t)$, $\chi_-(t)<x<\psi(t)$(或 $\psi(t)<x<\chi_+(t)$), 出发的极值后向特征线将在区间 $[x_-,\xi_-^t(0)]$(或 $[\xi_+^t(0),x_+]$)内被 $x$ 轴截获, 因而
\begin{align}
x-tf'(u(x\pm,t))=
\begin{cases}
x_-+o(t^{-1}), & \chi_-(t)<x<\psi(t),\\
x_++o(t^{-1}), & \psi(t)<x<\chi_+(t).
\end{cases}\label{align:ch11-11-7-5}
\end{align}
由于 $u(\chi_-(t),t)=u(\chi_+(t),t)=\bar u$,\cref{theorem:ch11-11-7-1} 蕴含 $u-\bar u=O(t^{-1})$, 故当 $t\to\infty$ 时, $f'(u)=f'(\bar u)+f''(\bar u)(u-\bar u)+O(t^{-2})$。此式连同\eqref{align:ch11-11-7-5}给出\eqref{align:ch11-11-7-4}。

最后, 将\eqref{align:ch11-11-7-4}中的 $u$ 代入\eqref{align:ch11-11-7-2}, 即得\eqref{align:ch11-11-7-3}。证毕。
\end{proof}

我们将用上述命题描述周期解的渐近行为。

\begin{theorem}\label{theorem:ch11-11-7-3}
当初值 $u_0$ 是周期的, 均值为 $\bar u$ 时, 则当 $t\to\infty$, 可容许解 $u$ 以速率 $o(t^{-1})$ 趋于一个周期锯齿剖面, 由形如\eqref{align:ch11-11-7-4}的小波构成。每个周期内小波(或齿)的个数等于每个周期内分界线的个数, 等价地, 等于 $x$ 轴上任意周期长度区间内函数 $u_0-\bar u$ 的原函数取极小值的点数。特别地, 在每个周期区间上 $u_0-\bar u$ 的原函数仅在单点取极小值的一般情形下, $u$ 趋于每个周期单齿的锯齿状剖面。
\end{theorem}

\begin{proof}
这是\cref{theorem:ch11-11-4-1}与\cref{theorem:ch11-11-7-2}的直接推论。若 $u_0$ 周期, \eqref{align:ch11-11-4-2}仅当 $\bar u$ 是 $u_0$ 的均值、且在 $u_0-\bar u$ 的原函数取极小值的点 $\bar x$ 处成立。这样的点集显然在周期平移下不变, 并在每个周期长度区间内至少含一个(一般恰为一个)点。
\end{proof}

\end{document}
